\section{HPS apparatus and data sets}
\label{sec:datasets}

\subsection{Detector configurations for the prompt search}

HPS is a forward spectrometer optimized for narrow $e^+e^-$ resonances and displaced
vertices in the high-rate environment immediately downstream of a thin tungsten target.
Its prompt $A^\prime \rightarrow e^+e^-$ sensitivity is set primarily by three
subsystems: the Silicon Vertex Tracker (SVT), which determines the mass and vertex
resolution; the electromagnetic calorimeter (ECal), which provides the trigger and
track--cluster matching; and the beamline/target geometry, which defines the forward
acceptance and the occupancy environment \cite{HPSExperiment2022,HPS2016PromptLong}.

The 2015 and 2016 engineering runs used the baseline six-station SVT. In that
configuration, six axial/stereo tracking stations were distributed between
$z=\SIrange{10}{90}{cm}$ downstream of the target, with the first active silicon only
about \SI{1.5}{mm} from the beam plane \cite{HPS2016BumpInternal,HPSExperiment2022}.
The ECal geometry used in the engineering runs was already close to the final HPS
layout: two PbWO$_4$ crystal halves with the highest-rate crystals nearest the beam
removed, leaving the beam to pass between the two halves
\cite{HPS2016BumpInternal,HPSExperiment2022}.

Before the first physics run in 2019, HPS upgraded the front tracker. The 2019 and
2021 configurations added a seventh double layer at $z=\SI{5}{cm}$ and used thin,
slim-edge sensors with short split striplets
\cite{HPSExperiment2022,HPSSlimEdgeSensors2020}. Moving the first precision
measurements closer to the beam reduced material and occupancy while improving
acceptance for low-mass, small-opening-angle pairs. The ECal crystal layout was
unchanged, but a positron hodoscope was installed before the 2019 run to support a
single-arm trigger for events near the ECal beam gap \cite{HPSExperiment2022}.
Figure~\ref{fig:hps-svt-eras} compares the tracker layouts used in the engineering and
physics runs.

\begin{figure}[H]
\centering
\begin{subfigure}[t]{0.90\linewidth}
  \centering
  \graphicorplaceholder{\linewidth}{apparatus_figs/hps_2015_2016_svt_baseline.png}
  \caption{Baseline six-layer SVT used in the 2015 and 2016 engineering runs
  \cite{HPS2016BumpInternal}.}
\end{subfigure}

\vspace{0.9em}

\begin{subfigure}[t]{0.78\linewidth}
  \centering
  \graphicorplaceholder{\linewidth}{apparatus_figs/hps_2019_2021_svt_upgrade.png}
  \caption{Front-layer upgrade used in the 2019 and 2021 physics runs, including the
  added seventh double layer and slim-edge front sensors
  \cite{HPSExperiment2022,HPSSlimEdgeSensors2020}.}
\end{subfigure}
\caption{SVT configurations used in the prompt-resonance analyses. The 2015/2016
baseline tracker and the 2019/2021 tracker differ in both layer count and the design
of the innermost stations. The ECal crystal layout was unchanged; the physics-run
configuration added a positron hodoscope.}
\label{fig:hps-svt-eras}
\end{figure}

\subsection{Data-taking campaigns and analysis scope}

This note contains two related analysis states. The accepted v4.2 combined result uses
the full 2015 and 2016 engineering-run samples together with a distributed 10\% sample
from the 2021 physics run. The later v4.9.5 pass examines that 2021 sample alone with a
different lower support edge. The engineering data share the baseline tracker used in
the published prompt searches; the 2021 data were recorded with the upgraded tracker
and positron hodoscope. Although the 2019 run used the same upgraded detector, no 2019
data enter either analysis.

Table~\ref{tab:dataset-summary} summarizes the run periods, analyzed samples, search
intervals, and GP fit supports. A mass hypothesis is tested only inside the search
interval, whereas the GP is trained on the wider support after removing the local
$2.25\,\sigma_m$ exclusion.

\begin{table}[H]
\centering
\small
{\setlength{\tabcolsep}{4pt}
\begin{tabularx}{\linewidth}{@{} P{0.20\linewidth} P{0.12\linewidth} P{0.28\linewidth} P{0.16\linewidth} Y @{}}
\toprule
Run period & Beam energy & Sample & Search interval & v4.2 GP fit support \\
\midrule
2015 engineering run &
\SI{1.056}{GeV} &
\SI{1170}{nb^{-1}} &
\SIrange{19}{90}{MeV} &
\SIrange{14}{135}{MeV} \\
2016 engineering run &
\SI{2.3}{GeV} &
\SI{10608}{nb^{-1}} &
\SIrange{39}{180}{MeV} &
\SIrange{30}{210}{MeV} \\
2021 physics run &
\SI{3.74}{GeV} &
Native 10\% sample from a $\sim\SI{160}{pb^{-1}}$ parent run &
\SIrange{50}{250}{MeV} &
\SIrange{40}{300}{MeV} \\
\bottomrule
\end{tabularx}
}
\caption{Data sets and mass domains in the accepted v4.2 three-campaign
configuration. The 2015 and 2016 inputs are the full engineering-run samples; the
2021 input is the distributed native 10\% sample. The wider GP supports provide
training sidebands beyond both ends of each search interval. The 2016 lower edge is
\SI{30}{MeV} because a \SI{35}{MeV} edge leaves no rebinned training-bin center below
the exclusion at the \SI{39}{MeV} search endpoint. The corresponding published
prompt searches covered \SIrange{19}{81}{MeV} in 2015 and
\SIrange{39}{179}{MeV} in 2016 \cite{HPS2015DarkPhoton,HPS2016PromptLong}.}
\label{tab:dataset-summary}
\end{table}

The later v4.9.5 study changes only the lower support edge for a 2021-only extraction:
it uses 36--300~MeV while retaining the 50--250~MeV search interval. The accepted v4.2
three-campaign configuration in Table~\ref{tab:dataset-summary} continues to use
40--300~MeV for 2021.

Integrated luminosity and analyzed fraction do not enter the GP interpolation. They
enter later, with the target, acceptance, efficiency, and radiative fraction, when an
extracted prompt yield is converted to an equivalent $\epsilon^2$ limit.

The 2015 and 2016 samples are fully unblinded. The 2021 input is distributed across
the run rather than drawn from a contiguous block, so it samples varied run conditions
while leaving most of the exposure unopened. The ROOT file does not contain enough
exposure metadata to infer an integrated luminosity; ``10\%'' describes how the sample
was constructed. The three-campaign result combines the full 2015 and 2016 samples
with this 2021 sample and is not the final result for the 2021 run.

Figure~\ref{fig:dataset-mass-shape-overlay} compares the unit-normalized invariant-mass
shapes. Figure~\ref{fig:dataset-mass-distributions-log} shows the same spectra in
events per MeV and marks the v4.2 fit supports and search intervals.

\begin{figure}[H]
\centering
\graphicorplaceholder{0.86\linewidth}{dataset_summary_figs/invariant_mass_distributions_2015_2016_2021_normalized.png}
\caption{Unit-normalized invariant-mass shapes for the full 2015 and 2016 samples
and the native 2021 10\% sample. Each histogram is divided by its event count and bin
width, so the comparison shows spectral shape rather than sample size.}
\label{fig:dataset-mass-shape-overlay}
\end{figure}

\begin{figure}[H]
\centering
\graphicorplaceholder{0.98\linewidth}{dataset_summary_figs/invariant_mass_distributions_2015_2016_2021_log.png}
\caption{Invariant-mass distributions for the three inputs in the accepted v4.2
three-campaign configuration. Gray shading marks the GP fit support and gold shading
marks the narrower interval in which signal hypotheses are tested. The mass-dependent
$2.25\,\sigma_m$ exclusion is removed from the GP support separately at every tested
mass.}
\label{fig:dataset-mass-distributions-log}
\end{figure}
